% ---
% Reemplaza líneas 2666–2719 de Docs/Technical_Report.tex.
% Decisiones que lo respaldan: D-006, D-008, D-009, D-023, H-007.
% ---

\subsection{Módulo 3: Extracción de características}

El módulo de extracción transforma cada par (imagen, máscara) producido por M2 en un vector
numérico $\mathbf{x}\in\mathbb{R}^{67}$ que describe la morfología de la lesión. Se emplea
PyRadiomics \cite{vanGriethuysen2017} sobre la imagen original, sin filtros, con cuatro familias de
características:

\begin{itemize}
  \item \textit{Primer orden} (18): estadísticas de la distribución de intensidades dentro de la
        lesión, como la media, los percentiles, la energía, la asimetría y la curtosis.
  \item \textit{Forma 2D} (9): descriptores geométricos del contorno, como el perímetro, los
        ejes mayor y menor, el diámetro máximo, la elongación y la esfericidad.
  \item \textit{GLCM} (24): textura a partir de la co-ocurrencia de pares de niveles de gris.
  \item \textit{GLRLM} (16): textura a partir de la longitud de las secuencias de píxeles con el
        mismo nivel de gris.
\end{itemize}

Las intensidades se discretizan en un número fijo de 32 niveles (\texttt{binCount}), no en un
ancho fijo de nivel. El número fijo es lo recomendado para modalidades de intensidad arbitraria,
como la película digitalizada del CBIS-DDSM, cuyos valores no tienen una escala física. Con un ancho
fijo, los datos de 16 bits generaban más de mil niveles. Las matrices de textura usan distancia 1 y
promedian las cuatro direcciones del plano. Como las mamografías no traen espaciado de píxel, las
características de forma se expresan en píxeles.

La extracción se aplica a las 3\,568 lesiones y a la variante con CLAHE de M2, sin fallos ni
valores faltantes. Masas y calcificaciones se extraen y se guardan por separado.

Como alternativa se contempló un codificador convolucional preentrenado
(ResNet-18 o EfficientNet-B0). Se eligió PyRadiomics porque cada característica tiene un significado morfológico
definido, lo que permite interpretar qué propiedades de la lesión conserva o pierde cada
transformación.
